\documentclass[12pt]{article}
\usepackage[margin=1in]{geometry}
\usepackage{enumitem}
\usepackage{titlesec}
\usepackage{parskip}
\usepackage{graphicx}
\usepackage{xcolor}
\usepackage{hyperref}
\usepackage{fontenc}

\title{\textbf{Appleby et al 2020 RG Notes}\\
\large Appleby, Joyce, Elizabeth Covington, David Hoyt, Michael Latham, Allison Sneider, Joyce Appleby, Elizabeth Covington, David Hoyt, Michael Latham, and Allison Sneider. ``Clifford Geertz (1973): Thick Description: Toward an Interpretive Theory of Culture'' [2020].\\
\textit{Knowledge and Postmodernism in Historical Perspective}: 310-323.}
\author{}
\date{}

\begin{document}

\maketitle

``If you want to understand what a science is, you should look in the first instance not at its theories or its findings, and certainly not at what its apologists say about it; you should look at what the practitioners of it do.'' (Appleby et al., 2020, p. 1)

Core argument: phenomena that may be observed as identical may actually be deeply different, both in cause, meaning, practice, etc. -- that is, just because we observe two identical outcomes does not suggest that we can characterize the two processes as identical

Particularly applied to culture:

\end{document}
